\documentclass[11pt,letterpaper]{article}
\usepackage[T1]{fontenc}
\usepackage{lmodern,microtype}
\usepackage[margin=1in]{geometry}
\linespread{1.0}
\setlength{\parskip}{0pt}
\setlength{\parindent}{1.5em}
\usepackage{amsmath,amssymb,amsthm,graphicx,xcolor,tikz,booktabs,tabularx,array,longtable}
\usepackage{listings}
\lstset{basicstyle=\ttfamily\footnotesize,breaklines=true,columns=fullflexible}
\usepackage[hidelinks]{hyperref}
\usepackage{cleveref}
\newtheorem{theorem}{Theorem}[section]
\newtheorem{proposition}[theorem]{Proposition}
\newtheorem{lemma}[theorem]{Lemma}
\newtheorem{corollary}[theorem]{Corollary}
\theoremstyle{definition}
\newtheorem{definition}[theorem]{Definition}
\newtheorem{example}[theorem]{Example}
\newtheorem{conjecture}[theorem]{Conjecture}
\theoremstyle{remark}
\newtheorem{remark}[theorem]{Remark}
\author{Matthew Long\\\small The YonedaAI Collaboration, YonedaAI Research Collective\\\small Chicago, IL\\\small\texttt{matthew@yonedaai.com}}
\date{26 November 2026}
\newcommand{\series}[1]{\\[4pt]{\normalsize #1 of \emph{Categorical Null-Cone Quantum Gravity}}}
\newcommand{\papersidebar}[1]{\AddToHookNext{shipout/foreground}{\begin{tikzpicture}[remember picture,overlay]\node[rotate=90,anchor=south,font=\footnotesize\sffamily\color{gray},inner sep=0pt] at ([xshift=30pt,yshift=.5\paperheight]current page.south west) {GrokRxiv:2026.11.#1\quad [hep-th]\quad 26 Nov 2026};\end{tikzpicture}}}

\title{Derived Gauge Reduction and Quantum Observable Algebras\series{Part II}}
\newcommand{\cA}{\mathcal A}
\newcommand{\cN}{\mathcal N}
\newcommand{\cZ}{\mathcal Z}
\newcommand{\CC}{\mathbb C}
\newcommand{\RR}{\mathbb R}
\begin{document}
\ifdefined\HCode\else\papersidebar{reduction}\fi
\maketitle
\begin{abstract}
A boundary frame can carry a measurable charge even when a simultaneous change of two frames
is gauge. We compute the reduced Poisson algebra, the Koszul complex and a quantum normalizer
quotient for a family of linear frame constraints. A relative coordinate and its momentum
survive both classical and quantum reduction. Repeating a matching condition creates a class
in negative degree without changing either ordinary observable algebra. On a graph, a consistent
quantum level requires the signed charge offsets to sum to zero around every cycle. A family
whose constraint loses rank shows that specializing a reduced algebra can miss observables of
the exceptional fiber. We use displacement probes to compare corresponding Gaussian states
and exhibit a composition phase invisible in every observable automorphism. For the null-frame
boost integral we compute a regulated matching kernel and distinguish its continuum limit
from projection onto equal discrete eigenvalues. These calculations give explicit criteria
for comparing the area charge, preparation and gauge data attached to a causal diamond.
\end{abstract}

\section{The choice at a corner}
\label{sec:corner}
Suppose two observers describe the same cut using different normal frames. If each observer is
free to change the common reference, only a relative reading can be measured. If one observer
changes a frame while the physical field and the reference on the other side stay fixed, the
change can act on a charged state. A prescription that calls both transformations gauge has
already decided which observables can exist.

Chandrasekaran gives this distinction a geometric expression in his null-cone construction
\cite{nc}. His equation (3.40) varies a pullback by varying the field and the boundary
embedding. The corner term in (3.41) contains the boost parameter and the area form. Equation
(3.43) identifies the integrated boost generator with the corner area divided by $8\pi G_N$.
At fixed area, a boost acts on a state by a phase in the representation of \S4.2. A correct
dictionary for the algebra must respect this action whenever the boundary conditions retain
it. Donnelly and Freidel describe the general distinction between gauge transformations and
physical surface symmetries in extended phase space \cite{df}.

Long's site of diamonds specifies where local data can be compared \cite[Example 4.8]{long}.
His later finite reconstruction starts from relational data on a graph \cite{rep}. Neither
construction specifies a quantum moment map for changing a null frame. Conversely, a reduced
phase space for a particular null region does not specify which categorical probes distinguish
two boundary states. The proposed comparison poses a concrete question. Which map of
gravitational data carries two equivalent representatives to equivalent representatives and
carries a measurable boundary quantity to the same quantity in the enlarged region?

We study that question in polynomial mechanics. A boundary has a coordinate recording a frame
and a conjugate momentum recording its charge. The diagonal shift of two frames provides one
gauge direction. A third observer can check the matching using the difference of the two
readings. All rings used for the explicit calculations have characteristic zero. Each model
has finitely many canonical pairs and an infinite-dimensional polynomial algebra. No choice of
spacetime, smearing kernel or Newton constant enters these calculations.

\section{Descent and admissible observables}
Write $\pi:X\to X/{\sim}$ for the quotient by an equivalence relation. For any set $V$, a
function $a:X\to V$ has a unique factorization through $\pi$ exactly when $a(x)=a(x')$ for
every related pair $x\sim x'$. Indeed, the required factor must take $[x]$ to $a(x)$, and the
displayed equality is precisely the condition for this formula to have a value independent of
the representative. This argument applies to completely general equivalence relations.
Differential or operator structure on the quotient needs additional hypotheses.

If $i:X\to Y$ is equivariant for two actions of the same group $G$, the pullback of a
$G$-invariant function on $Y$ is invariant on $X$. The pullback satisfies
\[
 (i^*a)(gx)=a(i(gx))=a(gi(x))=a(i(x)).
\]
One can also transport along a group homomorphism by taking $i(gx)=\rho(g)i(x)$ and requiring
invariance under its image. The direction of this map matters. Restricting classical fields on
a large region to a smaller one can induce pullback of observables from the small region to
the large one. Constructing the restriction itself requires compatible boundary conditions. An
inclusion of point sets does not define it.

Let $R$ be a commutative Poisson algebra, $I$ an ideal and $\{I,I\}\subset I$. Define
\[
 \cN_{\mathrm{cl}}(I)=\{a\in R:\{a,I\}\subset I\}.
\]
Then $I\subset\cN_{\mathrm{cl}}(I)$, and $\cN_{\mathrm{cl}}(I)/I$ inherits a Poisson bracket.
The Leibniz and Jacobi identities prove closure under multiplication and bracket. Replacing
one observable by a multiple of a constraint changes its bracket with another observable by an
element of $I$. Representatives outside the normalizer can still define functions on a
constraint surface as points; their brackets may depend on the choice of representative. For
constraints generating a Hamiltonian action, reduction also asks that the restricted
observable be constant on the gauge orbits. In the following examples the two descriptions
agree explicitly.

\section{A pair of boundary frames}
\label{sec:pair}
Let $(q_L,p_L)$ and $(q_R,p_R)$ be the coordinates of $T^*\RR^2$, with
$\{q_j,p_k\}=\delta_{jk}$. Both momenta refer to the outward orientation on their respective
sides. Use the matching equation and frame transformation
\[
 n=p_L+p_R=0,\qquad (q_L,q_R)\longmapsto(q_L+t,q_R+t).
\]
The transformation is generated by $n$ with the convention $X_n f=\{f,n\}$. Its zero level
sets the two outward charges equal and opposite. The sign depends on how the right frame is
oriented.

\begin{proposition}\label{prop:class}
For $R=\CC[q_L,q_R,p_L,p_R]$ and $I=(n)$, the classical reduced algebra is $\CC[r,k]$ with
$\{r,k\}=1$, where $r=q_L-q_R$ and $k=p_L$ on the matching surface. The induced symplectic
form on the real reduced space is $dr\wedge dk$ in the convention $\omega=\sum_j dq_j\wedge
dp_j$.
\end{proposition}
\begin{proof}
Set $s=q_R$. In the new coordinates the canonical one-form is
\[
 p_L\,dq_L+p_R\,dq_R=k\,dr+n\,ds .
\]
Thus $(r,k)$ and $(s,n)$ are independent canonical pairs, and $X_n=\partial_s$. At $n=0$ the
restricted symplectic form has kernel spanned by $\partial_s$. Each orbit meets the surface
$s=0$ once. An element of $R/(n)=\CC[r,k,s]$ is invariant under every real translation exactly
when it is independent of $s$. The same condition is $\partial_s a\in(n)$ before taking the
quotient, hence $\cN_{\mathrm{cl}}(I)/I=\CC[r,k]$. The bracket and form on the slice follow
from the canonical one-form.
\end{proof}

A measurement before reduction illustrates what the normalizer discards. The single reading
$q_L=s+r$ takes different values along the same orbit. Its difference from $q_R$ survives. The
charge $p_R$ also survives, since it is $-k$ at the fixed level. A separate reduction on each
side by the full translation of its frame, with each moment map fixed at zero, sets both
charges to zero. Each singly reduced phase space is a point. Combining these points loses the
canonical pair $(r,k)$ that the diagonal reduction retains. The choice of a common action has
measurable content in the model.

A subgroup that preserves a boundary condition can still act on the reduced observables. A
shift of the left reading alone induces $r\mapsto r+a$ and fixes $k$. It preserves the
matching surface and commutes with the diagonal shifts. Its generator on the quotient is $k$.
Declaring it gauge would impose a further reduction. This is the algebraic reason to identify
both the null directions of the restricted presymplectic form and the charged boundary
symmetries before constructing any algebra of a subregion.

\section{Regular and repeated constraints}
For a polynomial algebra $R$ and a list of functions $f=(f_1,\ldots,f_m)$, use the
cohomological convention for the Koszul complex
\[
 K(R;f)=R\otimes\Lambda(e_1,\ldots,e_m),\qquad |e_i|=-1,\qquad de_i=f_i .
\]
The differential vanishes on $R$ and satisfies the graded Leibniz rule. Its zeroth cohomology
is $R/(f_1,\ldots,f_m)$. It computes the derived zero fiber of the specified map to affine
$m$-space: resolve the origin there by its free coordinate Koszul resolution and pull that
resolution back to $R$. Two maps with different targets can have the same ordinary zero set
and different derived zero fibers. More general excess intersections have global geometric
obstructions; the calculation below uses global linear coordinates \cite{derived}.

\begin{proposition}\label{prop:excess}
The complex $K(R;n)$ for the surface in Proposition~\ref{prop:class} has cohomology $R/(n)$ in
degree zero and vanishing cohomology in negative degrees. The complex for the map given by
$(n,n)$ has
\[
 H^0=R/(n),\qquad H^{-1}=(R/(n))\,\varepsilon,\qquad
 H^j=0\quad(j<-1).
\]
Here $\varepsilon$ is represented by $e_2-e_1$.
\end{proposition}
\begin{proof}
Multiplication by the coordinate $n$ is injective, so the two-term complex $R e_1\to R$
resolves $R/(n)$. In the second complex set $u=e_1$ and $v=e_2-e_1$. Then $du=n$ and $dv=0$.
As a complex this is a copy of $K(R;n)$ together with its shift by the degree of $v$. The
class $[v]$ survives because a boundary in its summand has coefficient divisible by $n$.
\end{proof}

The negative-degree class here remembers a relation between two tests of the same charge. It
has no interpretation as a new radiative polarization. If a boundary problem presents one
equation twice merely as a bookkeeping convention, the intended derived fiber uses the
independent equation. Keeping the doubled target defines a different derived problem.
Equivalences of presentations in a derived formulation must come with quasi-isomorphisms. The
dimensions above rule one out between these two complexes.

The source of a stabilizer is a separate choice. For the single effective $\RR$-action in
Proposition~\ref{prop:class}, every orbit is free, including the orbits on the repeated zero
locus. One can instead declare a group $\RR^2$ to act by shifting both frames by $t_1+t_2$.
Its subgroup $\{(t,-t)\}$ fixes every point. That group action and the map with two identical
constraints arise together if the moment map is defined for the whole redundant Lie algebra.
The stabilizer and the Koszul class then have related dimensions but occupy different parts of
the reduction. A quotient by the effective group removes the ineffective subgroup. Changing
the derived target removes the redundant equation.

One can also have a stabilizer with independent constraint equations in other settings. For
example, quotient the one-pair reduced plane by the finite symmetry $(r,k)\mapsto(-r,-k)$,
without adding any Hamiltonian constraint. The origin has isotropy of order two. Its invariant
functions $u=r^2$, $v=rk$, and $w=k^2$ satisfy $uw=v^2$. To see that they generate, take any
monomial $r^ak^b$ with $a+b$ even and group factors into pairs. The quotient surface is
singular at the origin. Its finite stabilizer arises without repeating a single equation in
the Koszul calculation.

\section{Quantum reduction at a fixed level}
\label{sec:quantum}
Fix a real $\hbar\ne0$. Let $\cA$ be the polynomial Weyl algebra over $\CC$ generated by
$R_0,K,S,N$, with $[R_0,K]=[S,N]=i\hbar\,1$ and all commutators between the two pairs zero.
Uppercase letters indicate operators. The antilinear involution fixes these four generators.
Denote by $B$ the Weyl subalgebra generated by $R_0,K$. Classical variables can be recovered
as principal commutative symbols in the formal-parameter family described below.

The left ideal $L=\cA N$ contains the operators which visibly annihilate a vector satisfying
$N\psi=0$, on a common invariant domain when such vectors exist. Acting by a representative of
an observable must preserve the constraint on states. Define the right normalizer of the left
ideal by
\[
 \cN(L)=\{a\in\cA:L a\subset L\}.
\]
Since $L$ is a left ideal, $L\subset\cN(L)$ and it is a two-sided ideal inside $\cN(L)$. Thus
the quotient $\cN(L)/L$ has a product. The algebraic definition continues to make sense when
solutions to the state constraint are distributions.

\begin{theorem}\label{thm:weyl}
The inclusion $B\hookrightarrow\cN(L)$ induces an isomorphism of associative algebras
\[
 B\xrightarrow{\ \simeq\ }\cN(L)/L .
\]
The same quotient has a presentation as a star algebra by the centralizer of $N$ modulo the
ideal generated by $N$ inside that centralizer.
\end{theorem}
\begin{proof}
Normal ordering puts every element uniquely in the form
\[
 a=\sum_{j,\ell\geq0} b_{j\ell} S^j N^\ell,\qquad b_{j\ell}\in B ,
\]
with finitely many terms. The ordered basis follows by commuting $N$ past $S$ using the scalar
commutator, or by acting with multiplication and differentiation on polynomials in the two
coordinates. In the differential representation the highest nonvanishing derivative
coefficient detects a nonzero normally ordered operator in characteristic zero. Terms with
$\ell>0$ span $L$. Modulo $L$ there is a unique representative $a_0(S)=\sum_j b_{j0}S^j$.

It suffices to check $Na\in L$ in the definition of $\cN(L)$. Since $aN\in L$, this is
equivalent to $[N,a]\in L$. Its class is the class of $-i\hbar\,\partial_S a_0(S)$. The
derivative vanishes exactly on constant polynomials, giving surjectivity from $B$. Its
intersection with $L$ is zero by the ordered basis, giving injectivity.

The centralizer $\cZ(N)$ contains exactly the terms independent of $S$, by differentiating
their coefficients before taking any quotient. Thus $\cZ(N)=B[N]$ and $\cZ(N)\cap L=B[N]N$.
This ideal is star stable there, and evaluation at $N=0$ identifies the star quotient with
$B$. Its image in $\cN(L)/L$ is the associative isomorphism just proved.
\end{proof}

A two-sided ideal of the full $\cA$ containing $N$ contains $[S,N]=i\hbar\,1$. At the chosen
nonzero value of $\hbar$, the quotient by that ideal vanishes. Also $L$ itself fails star
stability since $SN\in L$ while $(SN)^*=NS=SN-i\hbar$ has nonzero constant class. The physical
star from the centralizer thus cannot be computed by taking an arbitrary normalizer
representative, applying the kinematical star and passing to its class. The proof specifies
which observables carry a well-defined adjoint.

Let $B_\hbar$ now denote the polynomial algebra over $\CC[\hbar]$ with the same relative pair
and relation $[R_0,K]=i\hbar$. Normal ordering gives a free module with basis $R_0^aK^b$. Its
specialization at $\hbar=0$ is the reduced Poisson polynomial algebra as a commutative
algebra, with bracket obtained from $[\,\cdot\,,\,\cdot\,]/(i\hbar)$ followed by
specialization. The coefficient of $\hbar$ satisfies the Leibniz rule and has value one on the
ordered generators. This proves the correspondence for all polynomials. In contrast,
specializing the kinematical commutator condition $[N,a]=0$ at $\hbar=0$ before dividing
removes the gauge derivative completely. The action or its divided derivation is part of the
classical data.

The fixed-parameter calculation avoids the general analytic questions of quantization and
reduction. Bordemann, Herbig and Waldmann derive hypotheses under which their quantum BRST
reduction deforms the full classical reduced function algebra, and exhibit examples in which
its degree-zero cohomology has different behavior \cite{brst}. The calculation here has a
global gauge coordinate, constant commutators and a polynomial ordering. In gravity one has to
establish which replacement for each assumption is available.

\section{Linear boundary actions}
\label{sec:linear}
The paired-frame calculation extends to several matching equations. Extra equations may depend
on the others. Their dependence becomes a consistency test when choosing the quantum value of
a charge.

Let $V$ and $U$ be real vector spaces of dimensions $d$ and $m$. A linear map $T:U\to V$ sends
gauge parameters to frame shifts. Its rank is $a$. On $T^*V$ let a vector $u\in U$ act by
$q\mapsto q+Tu$ with $p$ fixed. The moment map $\mu:T^*V\to U^*$ satisfies
\[
 \mu(q,p)(u)=p(Tu).
\]
Write $H=\operatorname{im}T$, $F=\ker T$. The action factors through the effective additive
group $H$. A character value $\lambda\in U^*$ determines the level $\mu=\lambda$. A nonzero
$\lambda$ is allowed even though the group is abelian. Physically it fixes what the matching
prescription is to count as zero total charge.

Let $\cA_\hbar(V)$ be the polynomial Weyl algebra with position $Q_\xi$ linear in $\xi\in V^*$
and momentum $P_v$ linear in $v\in V$. Take $[Q_\xi,P_v]=i\hbar\,\xi(v)1$, with commuting
positions and commuting momenta. All real generators are self-adjoint. Work over the complex
numbers with real $\lambda$ and fixed nonzero real $\hbar$. Define
\[
 C_\lambda(u)=P_{Tu}-\lambda(u)1,\qquad
 L_\lambda=\sum_{u\in U}\cA_\hbar(V) C_\lambda(u).
\]
Only a basis of $U$ is needed in the sum. The constraints commute. Their effective translation
action is free even when $F$ is nonzero, because a nonzero translation of $V$ fixes no point.

\begin{theorem}\label{thm:linear}
The level is nonempty exactly when $\lambda|_F=0$. For an inconsistent level
$L_\lambda=\cA_\hbar(V)$. For a consistent level, after a choice of complement to $H$ the
following descriptions hold.
\begin{enumerate}
\item The classical reduced phase space for the effective group is symplectomorphic to $T^*(V/H)$.
\item The cohomology of the Koszul complex of the $m$ component equations is the ring of the unreduced level tensored with $\Lambda^\bullet F_{\CC}$, placing $F_{\CC}$ in degree $-1$. Here $F_{\CC}$ is the complexification of $F$. Its $b$th exterior power contributes in degree $-b$ for $0\le b\le m-a$.
\item The quantum normalizer quotient $\cN(L_\lambda)/L_\lambda$ is a Weyl algebra on $d-a$ canonical pairs. Its star structure comes from the joint centralizer of the constrained momenta.
\end{enumerate}
\end{theorem}
\begin{proof}
A point on the level must satisfy $p\circ T=\lambda$, so $\lambda$ must vanish on $F$.
Conversely such a $\lambda$ defines a functional on $H$ by $Tu\mapsto\lambda(u)$, which
extends to $V$. For $u\in F$, the corresponding quantum constraint is $-\lambda(u)1$. If a
nonzero scalar occurs, the ideal contains $1$.

For a consistent level choose splittings $V=H\oplus W$ and $U=U_0\oplus F$. Use coordinates
$s_j$ on $H$ and $r_b$ on $W$, with momenta $n_j,k_b$ in the corresponding dual bases. The map
$T|_{U_0}$ is invertible onto $H$. In a basis adapted to this map the first $a$ constraints
fix the momenta $n_j$ at numbers $\lambda_j$, and the last $m-a$ vanish identically. The
symplectic form restricts to $\sum_b dr_b\wedge dk_b$. The quotient by translations in $s$
leaves the stated cotangent bundle.

The shifted coordinates $n_j-\lambda_j$ form a regular sequence. In fact the constraint
complex factors over $\CC$ into independent two-term coordinate complexes, each exact in
degree $-1$, a polynomial ring in the other variables, and an exterior algebra with zero
differential on the generators in $F_{\CC}$. The cohomology follows by tensoring over a field.
When taking invariants under the effective translations afterward, the polynomial dependence
on the $s_j$ drops out. This last observation concerns ordinary invariants of the displayed
cohomology. A general derived quotient by a group requires its own cohomological construction.

For the quantum algebra put every $N_j-\lambda_j$ at the far right of each normally ordered
monomial. The terms containing any such factor span $L_\lambda$. Modulo $L_\lambda$, each
class has a unique representative polynomial in $S_1,\ldots,S_a$ with coefficients in the
residual Weyl algebra. Normalizing all the constraints requires each of its partial
derivatives in $S_j$ to vanish. Thus just the constant coefficient remains. The centralizer
consists of residual Weyl polynomials together with polynomials in the commuting $N_j$.
Evaluation $N_j=\lambda_j$ on this centralizer is a star map because the levels are real.
\end{proof}

At $\lambda=0$ the quotient identification needs fewer choices. The momentum $p$ then belongs
to $\operatorname{Ann}(H)\simeq(V/H)^*$, and a gauge orbit of $q$ specifies a point of $V/H$.
At a nonzero consistent level the allowed momenta form an affine translate of
$\operatorname{Ann}(H)$. Choosing an extension of the character from $H$ to $V$ identifies
this affine space with a vector space. The theorem makes no preferred choice among such
extensions. A comparison of different diamonds could depend on this choice through its action
on momenta.

There is a useful direct form of the affine choice problem. Suppose the same boundary level is
expressed by two real-valued linear functionals on $V$, $p_*$ and $p_*'$. Both extend the
fixed character on $H$. Their difference $\theta$ annihilates $H$, hence defines a cotangent
vector on the quotient. The change of origin of the reduced momentum acts by $k\mapsto
k-\theta$. In the quantum Weyl algebra such a real shift preserves commutators and the
adjoint. A named observable interpreted as a fixed absolute momentum nevertheless changes
under the shift. When physical preparation data fix a preferred zero, the comparison must
carry that zero along with the ring operations.

\section{Infinitesimal gauge cochains}
\label{sec:cochains}
A derived zero fiber still carries the gauge action. For translations both operations can be
examined in one complex. The following calculation also shows what it costs to retain an
ineffective gauge generator. It uses polynomial Lie algebra cochains for an abelian Lie
algebra. Other choices, including equivariant topological cohomology, ask different
cohomological questions.

Take a consistent character and the adapted coordinates of Theorem~\ref{thm:linear}. Write $R$
for the ring of complex polynomials in $(r,k,s,n)$. Introduce a degree $-1$ generator for each
component of the moment map as above. Let the entire abelian Lie algebra $U$ act by its given
translation on position coordinates and trivially on momentum coordinates and Koszul
generators. Form the Chevalley and Eilenberg cochains of this action on the Koszul algebra.
The new degree $+1$ generators are linear functions on $U$. These generators anticommute with
the odd Koszul generators under the usual total-degree sign rule.

In a basis of $U_0$ mapped to the coordinate translations of $H$, write $e_j$ for the active
Koszul generators and $c_j$ for the ghosts dual to those translations. In a basis of $F$ write
$v_i$ for its Koszul generators and $w_i$ for the additional Lie algebra cochains. The full
differential $D$ satisfies
\begin{equation}
 D e_j=n_j-\lambda_j,\qquad D s_j=c_j,\qquad
 Dv_i=Dw_i=Dc_j=Dn_j=0 .
 \label{eq:BRST}
\end{equation}
It vanishes on $r,k$. For homogeneous $x$ of total degree $|x|$ (Koszul degree plus Lie
cochain degree), its value on a product is $D(xy)=(Dx)y+(-1)^{|x|}x(Dy)$.
The two component differentials are
compatible because the constraints are invariant under translations. In particular $D^2=0$ on
each generator. Nilpotence alone does not identify the cohomology with the algebra of ordinary
observables.

\begin{proposition}\label{prop:ce}
The cohomology of the polynomial complex \eqref{eq:BRST} is
\[
 \CC[r,k]\otimes
 \Lambda^\bullet F_{\CC}\otimes
 \Lambda^\bullet F_{\CC}^*,
\]
with degrees $-1$ and $+1$ on the two exterior generating spaces, respectively. Each term
$\Lambda^i F_{\CC}\otimes\Lambda^j F_{\CC}^*$ contributes in total degree $j-i$. With
independent constraints and the effective group, this cohomology is concentrated in degree
zero and equals the ordinary reduced observable algebra.
\end{proposition}
\begin{proof}
For each $j$ the subcomplex with variables $n_j-\lambda_j,e_j$ is the two-term coordinate
resolution in Proposition~\ref{prop:excess}. The other factor with variables $s_j,c_j$ is the
polynomial complex
\[
 \CC[s_j]\ \xrightarrow{\ \partial_{s_j}\ }\ \CC[s_j]c_j .
\]
Its kernel consists of constants. In characteristic zero every polynomial coefficient has a
polynomial antiderivative, so the derivative is onto. Both factors have just $\CC$ in degree
zero as cohomology. Taking tensor products over a field gives the assertion, since the
residual polynomial variables, $v_i$ and $w_i$ all have zero differential.
\end{proof}

When $\dim F=1$ there is a class $v$ in degree $-1$ from the redundant equation and a class
$w$ in degree $+1$ from the ineffective Lie algebra generator. The product $vw$ is a nonzero
class of total degree zero with $(vw)^2=0$. Thus even degree-zero cohomology can remember a
pair of extra choices in this complex. Dropping the ineffective Lie algebra and keeping the
doubled equation deletes $w$ but retains $v$. Replacing both inputs by the effective ones
deletes both classes. The reduced ordinary observable algebras in Theorem~\ref{thm:linear} are
the ones obtained by ordinary effective reduction. This distinction prevents an inference from
a correctly computed cohomology group to a statement about the physical observables it should
represent.

The proof uses an explicit primitive in the orbit coordinate. That primitive need not satisfy
the boundary conditions of another problem. For comparison, replace one noncompact coordinate
$s$ by an angular coordinate and act on smooth periodic functions. The derivative of a
periodic function has integral zero around the circle, and any smooth zero-integral function
has a periodic primitive. The constant coefficient of a one-form therefore gives a
one-dimensional cokernel. The orbit action is still free. This degree-one class reflects the
cochain complex chosen on a compact orbit. It supplies a second reason to specify the
coefficient space and the treatment of boundary conditions in a gravitational reduction.

\section{A constraint of changing rank}
\label{sec:rank}
The rank hypothesis of Theorem~\ref{thm:linear} can fail along a family of boundary data. The
simplest example uses the same coordinates $(r,k,s,n)$ as the paired-frame model. Introduce a
real parameter $\tau$ which the Poisson bracket holds fixed. A parameter $u\in\RR$ acts on $s$
by $s\mapsto s+\tau u$. Its moment map is
\begin{equation}
 f_\tau=\tau n .
 \label{eq:rank-jump}
\end{equation}
For each $\tau\ne0$ the ordinary zero fiber has $n=0$, and translating $s$ removes a frame
coordinate. At $\tau=0$ every point satisfies the equation and the effective action is trivial.
Reduction of this individual fiber leaves two canonical pairs. Extending results for the
regular fibers to the exceptional value would miss a whole canonical pair.

\subsection{The relative zero fiber}
Over the polynomial parameter ring take $D=\CC[\tau,r,k,s,n]$. The family of ordinary zero
fibers has coordinate ring $D/(\tau n)$. This ring has $\tau$-torsion: the nonzero class of $n$
is killed by $\tau$. Thus this ordinary family is not flat over $\CC[\tau]$. There is a simple
geometric explanation. Its zero set has a component $n=0$ over every value of the parameter,
and an additional component, with $n$ arbitrary, lying entirely over $\tau=0$.

The equation $\tau n$ is a nonzero divisor in $D$. Its Koszul complex
\[
 D e \xrightarrow{\ \tau n\ }D
\]
therefore has no negative-degree cohomology over the whole family. After setting $\tau=0$ in
the complex, its differential vanishes. The exceptional fiber has a copy of $\CC[r,k,s,n]$ in
degrees $0$ and $-1$. This calculation uses specialization of the complex of free
$\CC[\tau]$-modules. In particular there is more information in that specialization than in
tensoring the cohomology groups with $\CC[\tau]/(\tau)$ after computing them on the total
space.

One can see the additional term without introducing Koszul generators. Use the two-term
resolution of $\CC[\tau]/(\tau)$ whose differential is multiplication by $\tau$. Tensor it with
$D/(\tau n)$. The kernel of the resulting differential is the ideal $(n)$ in $D/(\tau n)$.
Multiplication by $n$ identifies $\CC[r,k,s,n]$ with that kernel, because every term containing
both $\tau$ and $n$ vanishes and every surviving multiple of $n$ has a unique polynomial
coefficient at $\tau=0$. This identifies the first Tor contribution with the degree $-1$ class
in the specialized Koszul model. Thus the surplus class can be traced either to a dependent
equation in an individual fiber or to failure of flatness in a family.

The quantum discrepancy can be computed without any homological language. Keep $\hbar$ real and
nonzero and set $A_0=\cA_\hbar(\RR^2)$ in the variables of Theorem~\ref{thm:weyl}. For
polynomials in $\tau$ let
\[
 A=A_0[\tau],\qquad L=A\tau N,\qquad
 M=\{a\in A_0:A_0N a\subset A_0N\}.
\]
The rank-changing normalizer is defined using the single family of constraints. Reducing
instead after specializing the kinematical algebra uses the individual constraint $f_0=0$.

\begin{proposition}\label{prop:basechange}
The specialized map from the family normalizer quotient to the physical quotient of the
exceptional fiber is the strict inclusion
\[
 \bigl(\cN(L)/L\bigr)\big/\tau\bigl(\cN(L)/L\bigr)
 \ \simeq\ M\ \hookrightarrow\ A_0.
\]
For $\tau_0\ne0$, specialization at $\tau_0$ gives the one-pair Weyl algebra $B$.
\end{proposition}
\begin{proof}
It suffices to test the generator $\tau N$ in the normalizer condition. Since $\tau$ is central
and is a nonzero divisor in $A$, the relation $\tau N a\in \tau A N$ is equivalent to $N a\in A
N$. The latter holds coefficient by coefficient in $\tau$. Hence $\cN(L)=M[\tau]$. The
ordered-monomial calculation in Theorem~\ref{thm:weyl} gives $M=B+A_0 N$ as vector subspaces,
so $A_0 N\subset M$. It follows that $L\subset \tau M[\tau]$. The special fiber of $M[\tau]/L$
is $M$.

The constraint in the exceptional kinematical fiber generates the zero ideal of $A_0$. Its
normalizer is all of $A_0$. The induced map is the inclusion $M\to A_0$, which misses $S$
because $NS=SN-i\hbar$ lies outside $A_0N$. At $\tau=\tau_0\ne0$ the image of $L$ is $A_0N$
since $\tau_0$ is invertible, giving $M/A_0N=B$.
\end{proof}

Two observations make the strict inclusion visible on states. In a regular fiber the
wavefunctions solving the translation constraint have a fixed $N$-value and a free $R_0$
argument. At the exceptional fiber a normalized kinematical packet can depend on both $r$ and
$s$. Multiplication by $S$ distinguishes two packets translated in $s$. There is no family of
polynomial normalizer representatives whose value at zero is this operator. Also a
representative proportional to $N$ can survive at $\tau=0$ in the source algebra $M$, even
though at a regular level it vanishes in the quotient. Both effects follow from performing
different operations in a different order. The model requires no limit of Hilbert spaces as
$\tau$ approaches zero.

A regular family of null diamonds likewise needs a stated rank condition if its reduced phase
spaces are to have a common description. Equation~\eqref{eq:rank-jump} is an example of what to
test when that condition fails. It gives no basis for assigning this singularity to a caustic
or marginal surface in a spacetime solution. Such an assignment would require the linearized
constraints of that particular solution and their boundary conditions.

\section{Matching around a cycle}
\label{sec:cycle}
A number assigned to each cut may look consistent when two neighboring pieces are examined and
fail when the cuts close into a loop. The linear theorem turns this into a calculation.

Let $G=(V_G,E_G)$ be a finite oriented graph with $c$ connected components. Retain isolated
vertices. Put one frame position and one momentum at each vertex and assign a translation
parameter to each edge. The column of $T:\RR^{E_G}\to\RR^{V_G}$ for the edge $e:v\to w$ is the
vector $e_w-e_v$. The corresponding constraint fixes $p_w-p_v=\lambda_e$. Reversing the edge
reverses the sign of its column and of $\lambda_e$. This is a charge-difference convention;
the two-boundary convention of Section~\ref{sec:pair} reverses the momentum on one side.

\begin{corollary}\label{cor:graph}
The vertex matching problem has a nonzero quantum quotient exactly when the signed sum of
$\lambda_e$ around every cycle vanishes. In that case the quotient has $c$ canonical pairs.
The unreduced derived level has degree $-1$ Koszul cohomology free over the ordinary level
ring of rank $b=|E_G|-|V_G|+c$. Its successive exterior powers fill degrees $-j$ for $0\le j\le b$.
\end{corollary}
\begin{proof}
The transpose of the incidence matrix maps a vector of vertex values $p$ to its edge
differences. Such differences have vanishing sum around every cycle. Conversely choose one
root in each component. Set a vertex value to the signed sum of the edge assignments along any
path from the root. Vanishing cycle sums give path independence. A constant per component
remains free, hence $\dim\ker T^*=c$ and $\operatorname{rank}T=|V_G|-c$. Apply
Theorem~\ref{thm:linear}.
\end{proof}

In a triangle with arrows $1\to2\to3\to1$, fix $\lambda_{12}+\lambda_{23}+\lambda_{31}=0$. One
may solve the constraints by $p_1=t$, $p_2=t+\lambda_{12}$ and
$p_3=t+\lambda_{12}+\lambda_{23}$. The sum $Q=q_1+q_2+q_3$ is invariant under all edge shifts.
Its conjugate can be taken to be $K=(p_1+p_2+p_3)/3$. On this level
$K=t+(2\lambda_{12}+\lambda_{23})/3$, so it differs from the free momentum $t$ by a fixed offset. Indeed $\{Q,K\}=1$ and the quantized
operators satisfy $[Q,K]=i\hbar$. The Koszul class represented by $e_{12}+e_{23}+e_{31}$ is
closed. It has degree $-1$ and spans the excess part over the ring of the unreduced level.
Thus the classical observable pair, the quantum observable pair and the redundant matching
equation can each be named in the same example.

Now put $\lambda_{12}=\lambda_{23}=0$ and $\lambda_{31}=\hbar\alpha$ with real $\alpha\ne0$.
Every proper two-edge subgraph of the triangle has consistent levels. Adding the third edge
gives
\[
 C_\lambda(e_{12})+C_\lambda(e_{23})+C_\lambda(e_{31})
   =-\hbar\alpha\,1 .
\]
The full left ideal contains a nonzero scalar and the quotient vanishes. At $\hbar=0$ all
three level assignments reduce to the ordinary unshifted matching. This is an obstruction to
lifting those particular levels together, despite consistency on each tree. It is an algebraic
test of an assignment of charge offsets. Calling it a gravitational anomaly would require an
invariant construction of the offsets in gravity and proof that no allowed redefinition
removes it.

A change of vertex momentum origins $p_v\mapsto p_v+b_v$ changes the edge character by
$b_w-b_v$. The signed sum around any cycle stays fixed. Conversely Corollary~\ref{cor:graph}
identifies all removable assignments as vertex differences. A spanning tree gives an
economical measurement: calibrate vertex origins along the tree, then measure each remaining
edge once. If an edge disagrees, no new vertex calibration can satisfy that edge while
retaining every tree match. In the formal-parameter algebra a nonzero scalar multiple of
$\hbar$ implies $\hbar$-torsion rather than instant collapse. Inverting $\hbar$ or
specializing to a nonzero numerical value gives the zero quotient of the example.

The dimension $c$ counts components of this chosen graph of gauge translations. It gives no
degree-of-freedom count for a gravitational partition of spacetime. The assumptions making the
count possible are that each frame variable is one-dimensional, translations commute,
constraints are linear and their only dependencies are graph cycles. A nonabelian charge
matching problem instead needs a specified composition law for changes of frame, and
representation data can distinguish sectors with the same classical label.

\section{Transport of a constraint ideal}
\label{sec:transport}
The arrows of a putative diamond category must do more than relate the unconstrained algebras.
For algebras $\cA_j$ with left ideals $L_j$, $j=1,2$, suppose a unital algebra map
$\alpha:\cA_1\to\cA_2$ satisfies
\begin{align}
 \alpha(L_1)&\subset L_2, \label{eq:ideal}\\
 \alpha\bigl(\cN(L_1)\bigr)&\subset\cN(L_2). \label{eq:norm}
\end{align}
Then $[a]\mapsto[\alpha(a)]$ induces a unital algebra map of normalizer quotients.
Equation~\eqref{eq:ideal} gives independence of representatives. Equation~\eqref{eq:norm}
ensures the image represents an observable. These conditions compose. If each physical adjoint
comes from a star-stable centralizer presentation, a sufficient additional requirement is that
$\alpha$ act as a star map on the centralizer and preserve its chosen target centralizer. An
embedding of unconstrained algebras by itself establishes none of the conditions concerning
target constraints.

A stronger hypothesis reduces the number of separate checks. Suppose $\alpha(L_1)\subset L_2$
and the left ideal generated by $\alpha(L_1)$ in $\cA_2$ equals $L_2$. If $a\in\cN(L_1)$ and
$b=\sum_i c_i\alpha(\ell_i)\in L_2$, then
\[
 b\,\alpha(a)=\sum_i c_i\alpha(\ell_i a)\in L_2 .
\]
Here multiplicativity of $\alpha$ allows each product to pass inside, and $\ell_i a\in L_1$
because $a$ normalizes the left ideal. Thus equation~\eqref{eq:norm} holds. An enlarged boundary can acquire independent constraints,
in which case this sufficient hypothesis may fail. The original normalizer test still makes
sense.

For a simple obstruction let $\alpha_t$ act on the two-pair algebra by $N\mapsto N+t\hbar$ and
fix $R_0,K,S$, with $t\in\RR$. It preserves the star and all commutators. This formula also
defines a map over $\CC[\hbar]$, and its specialization at $\hbar=0$ is the identity on all
four generators. If both regions impose $N=0$ at the quantum level, $\alpha_t(N)$ lies in $\cA
N$ exactly when $t=0$. A reduced representative and the same representative plus $N$ then have
images differing by $t\hbar$ modulo the target ideal. For $t\ne0$, equation~\eqref{eq:ideal}
fails and there is no map induced by this rule on all representatives.

The abstract physical algebras in the shifted example are nevertheless isomorphic. They have
common representatives in $B$, on which $\alpha_t$ is the identity. One can either use that
particular section to define a different prescription, or change the target constraint to
$N+t\hbar=0$, in which case the original map does respect the chosen levels. The distinction
is an experimentally meaningful one when levels and observables have fixed operational
definitions. A null-cone correspondence must specify the frame charge detected by a
preparation before treating changes of origin as equivalent.

\section{Physical states at a noncompact boundary}
An algebraic quotient does not automatically furnish a subspace of normalized physical
vectors. In the real-coordinate representation on $L^2(\RR^2,dr\,ds)$ take initially the
Schwartz domain and let
\[
 R_0=r,\qquad K=-i\hbar\partial_r,\qquad S=s,\qquad N=-i\hbar\partial_s .
\]
The constraint would set a state independent of $s$. Any such function in $L^2(\RR^2)$ has
zero norm or infinite norm in the $s$-direction, so the zero eigenspace of the self-adjoint
momentum $N$ contains only the zero vector. The reduced Weyl algebra in Theorem~\ref{thm:weyl}
still has its usual Schr\"odinger representation on $L^2(\RR,dr)$. One must account for the
gauge orbit when comparing the state spaces.

For Schwartz wavefunctions, consider the map
\[
 (\eta\psi)(r)=\int_{\RR}\psi(r,s)\,ds .
\]
Its normalization fixes a choice of measure on the gauge group. Integration by parts gives
$\eta(N\psi)=0$, and $R_0,K$ commute through this averaging map on the common domain. Every
Schwartz function $f(r)$ has a preimage $f(r)h(s)$ with $h$ Schwartz and $\int h(s)\,ds=1$.
The positive semidefinite form
\[
 (\psi,\phi)_{\mathrm{red}}
   =\int_{\RR}\overline{(\eta\psi)(r)}(\eta\phi)(r)\,dr
\]
therefore has quotient completion $L^2(\RR,dr)$. This identifies the physical scalar product
for the model by a construction that can be checked on representatives. It does not turn the
average into a bounded projection on the kinematical Hilbert space.

For example choose a Schwartz function $h$ with nonzero integral and set
$h_b(s)=b^{-1/2}h(s/b)$ for $b>0$. Its kinematical norm is constant while $\int h_b=\sqrt
b\,\int h$. For a fixed nonzero $f(r)$, the norms of $\eta(fh_b)$ tend to infinity. No bounded
extension of $\eta$ from the full kinematical Hilbert space can have this action on the
Schwartz domain. Replacing the gauge average by an ordinary Hilbert projection would give a
different answer.

A real level $\lambda$ changes the test functional to
\[
 (\eta_\lambda\psi)(r)=\int_{\RR} e^{-i\lambda s/\hbar}\psi(r,s)\,ds .
\]
Then $\eta_\lambda((N-\lambda)\psi)=0$ by integration by parts. A multiplication operator
$U_t\psi=e^{its}\psi$ satisfies $N U_t=U_t(N+t\hbar)$ and
$\eta_{\lambda+t\hbar}(U_t\psi)=\eta_\lambda(\psi)$. The phase implements a change of level
provided the physical state prescription changes along with the operator. The quantum shift
that obstructs descent at two fixed equal levels now has a precise intertwiner between two
different levels.

These arguments separate three assertions about a proposed local quantum system. A polynomial
star algebra may admit a representation on a dense domain. A bounded algebra of observables
may be generated by exponentiating suitable self-adjoint operators. A von Neumann algebra
depends on the selected representation and its closure. A theorem about a star-algebra
quotient cannot by itself identify the latter two. Chandrasekaran specifies deformations of
the action and a positive state in constructing his representation \cite[\S\S4.1, 4.2]{nc}. An
independent categorical construction must identify the comparable operators, the state, and
any limiting operation used to form a closure.

\section{A state comparison on the relative pair}
\label{sec:states}
A comparison can preserve the constraint algebra and still assign different probabilities to a
boundary measurement. The relative pair admits a direct calculation of those probabilities.
Work in the Schr\"odinger representation of Section~\ref{sec:quantum}, after group averaging
has removed $s$. For real $a,b$ define a bounded unitary on $L^2(\RR,dr)$ by
\begin{equation}
 (W(a,b)f)(r)=
 \exp\bigl(ia(r+b/2)/\hbar\bigr)\,f(r+b).
 \label{eq:displacement}
\end{equation}
Here $a$ has units of momentum and $b$ has units of position. Translation preserves Lebesgue
measure, so $W(a,b)$ is indeed unitary. Direct multiplication gives
\[
 W(a,b)W(c,d)=
 \exp\bigl(-i(ad-bc)/(2\hbar)\bigr)W(a+c,b+d).
\]
It follows that $W(a,b)^*=W(-a,-b)$. These are bounded probes of the pair. Their product is
meaningful for every Hilbert vector, in contrast with a high polynomial moment of an unbounded
generator. The symmetrically ordered characteristic function of displacement operators also
appears in the phase-space treatment of Cahill and Glauber \cite{cahill}.

For a unit vector $\psi$, write $\chi_\psi(a,b)=\langle\psi,W(a,b)\psi\rangle$. An algebra of
measured observables may be defined as the operator norm closure of the finite linear span of
these $W(a,b)$ in this representation. Equality of two states on all the displacements implies
equality on that algebra by linearity and norm continuity. For the calculation it suffices to
work with normalized vectors and the finite linear span.

\begin{proposition}\label{prop:gaussian}
Let $\sigma>0$ and $x,p\in\RR$. The normalized packet
\[
 g_{x,p,\sigma}(r)=
 (2\pi\sigma^2)^{-1/4}
 \exp\left(-\frac{(r-x)^2}{4\sigma^2}
           +\frac{ip(r-x)}{\hbar}\right)
\]
has characteristic function
\begin{equation}
 \chi_{x,p,\sigma}(a,b)=
 \exp\left(\frac{i(ax+bp)}{\hbar}
           -\frac{\sigma^2 a^2}{2\hbar^2}
           -\frac{b^2}{8\sigma^2}\right).
 \label{eq:character}
\end{equation}
Equality for all $a,b$ distinguishes the three parameters within this family, up to the
irrelevant common phase of a wavefunction.
\end{proposition}
\begin{proof}
The squared modulus is the normal probability density with mean $x$ and variance $\sigma^2$. In
the integral for $\chi$ the product of the two wavefunctions has the factor $e^{ipb/\hbar}$.
Put $u=r-x+b/2$. Completing the square gives
\[
 (r-x)^2+(r+b-x)^2=2u^2+b^2/2 .
\]
The remaining phase in \eqref{eq:displacement} is $e^{ia(u+x)/\hbar}$. Integrating
$e^{iau/\hbar}$ against the centered normal density yields $e^{-\sigma^2 a^2/(2\hbar^2)}$,
proving \eqref{eq:character}. The first derivatives at the origin recover $x,p$; the modulus
along the $a$-axis determines $\sigma^2$. Positivity of the width fixes its square root.
\end{proof}

A chosen pair of preparations can now test an affine change of boundary reference. For displacements let
\[
 \alpha_{\delta,\epsilon}\bigl(W(a,b)\bigr)=
 e^{i(a\delta+b\epsilon)/\hbar}W(a,b).
\]
The scalar is a character of the additive parameter space. It respects products and adjoints.
The map is unitarily implemented in this representation by conjugation with a displacement, as
follows also from the product law. The underlying operator prescription corresponds to
$R_0\mapsto R_0+\delta$ and $K\mapsto K+\epsilon$.

Let the first observer use parameters $(x_1,p_1,\sigma_1)$ and the second use
$(x_2,p_2,\sigma_2)$. Equality of the first state with the pullback of the second then holds on
the displacement algebra exactly when
\begin{equation}
 x_1=x_2+\delta,\qquad p_1=p_2+\epsilon,\qquad \sigma_1=\sigma_2 .
 \label{eq:state-match}
\end{equation}
Indeed \eqref{eq:character} gives these conditions and proves their sufficiency. A mismatch of
means can be accounted for by an identified offset in the frame origin. A mismatch of widths
remains after such a translation. Two preparations with zero means and different widths agree
on the means of both generators and on the canonical commutator, yet give different values even
for $\operatorname{Re}\chi(a,0)$ at every $a\ne0$. Thus a second moment can separate
comparisons that a measurement of charge and position means cannot distinguish.

Rescaling the unit of the frame reading changes the test. For $\gamma>0$ take
$\alpha_\gamma(W(a,b))=W(\gamma a,b/\gamma)$. The product phase in \eqref{eq:displacement} is
unchanged, and dilation of the Schr\"odinger coordinate implements this star automorphism.
Pulling back the second Gaussian state gives parameters $(\gamma
x_2,p_2/\gamma,\gamma\sigma_2)$. To use the width as a test one must fix the physical
normalization of the probe or include its prescribed rescaling. This dependence is relevant for
null frames, where a convention for normalizing the null normals enters a comparison of
readings.

Agreement on one commutative family of probes can also fail to identify a state. The packets
$g_{x,p,\sigma}$ and $g_{x,p',\sigma}$ for distinct $p,p'$ have identical probability
distributions for $R_0$. Every bounded function of $R_0$ therefore has the same expectation in
the two states. A translation probe with $a=0$ has expectation $e^{ibp/\hbar-b^2/(8\sigma^2)}$
and distinguishes the states for any $b$ with $b(p-p')/\hbar$ outside $2\pi\mathbb Z$. In a
boundary measurement this distinction asks which noncommuting operations the two constructions
allow a preparation to undergo. Knowing every probability for the charge alone would still give
only the diagonal spectral data of that charge.

The measurements in this section can be made after gauge averaging on representatives. For a
normalized residual packet $g$ choose any Schwartz function $h$ with $\int h(s)\,ds=1$.
Applying the bounded displacement in the $r$ variable to $g(r)h(s)$ before applying $\eta$
produces $W(a,b)g$. Replacing $h$ by any other function with the same integral changes that
kinematical representative but leaves every residual displacement expectation the same. For a
nonzero level the corresponding condition is $\int e^{-i\lambda s/\hbar}h(s)\,ds=1$. Thus
equality of the probability tests is compatible with the level-dependent state prescription,
not just with a preferred gauge slice.

\section{Translations implemented up to phase}
\label{sec:phases}
An algebra map can respect products on the nose when the operators implementing a boundary
symmetry respect composition only up to a phase. This issue is independent of the level
obstruction in Corollary~\ref{cor:graph}. In that example the global quotient vanished. The
following example has a nonzero Hilbert space and an exact action on every bounded observable.

On the reduced Hilbert space define $U(a)f(r)=f(r-a)$ and $V(b)f(r)=e^{ibr/\hbar}f(r)$. Both
shift the surviving phase-space data. They are boundary operations on the surviving canonical
pair and do not impose two additional gauge conditions. Direct substitution gives
\begin{equation}
 U(a)V(b)=e^{-iab/\hbar} V(b)U(a).
 \label{eq:weylphase}
\end{equation}
The scalar disappears in the inner automorphisms $A\mapsto U(a)AU(a)^*$ and $A\mapsto
V(b)AV(b)^*$. These automorphisms commute even when the two unitaries anticommute. They
preserve the residual displacement algebra since conjugating a displacement multiplies it by a
character.

\begin{proposition}\label{prop:phase}
Fix two nonzero steps $a_0,b_0$. The two commuting automorphisms associated with these steps
define an action of $\mathbb Z^2$ on bounded operators. In the given representation there is a
pair of implementing unitaries giving a true representation of $\mathbb Z^2$ if and only if
$a_0b_0/(2\pi\hbar)$ is an integer.
\end{proposition}
\begin{proof}
In the algebra of all bounded operators on the Hilbert space, unitaries implementing the same
inner automorphism differ by a scalar. Indeed their quotient commutes with all bounded
operators, including every rank-one operator, so it is a multiple of the identity. Changing
scalar phases of $U(a_0),V(b_0)$ leaves their group commutator unchanged. Its value is
$e^{-ia_0b_0/\hbar}1$ by \eqref{eq:weylphase}. A representation of the free abelian group on
two generators requires that commutator to be one. When the condition holds, the assignment
$(m,n)\mapsto U(a_0)^m V(b_0)^n$ gives the representation.
\end{proof}

A way to compute the obstruction without considering all elements of the group is to compare
the two paths around the elementary square. For two intertwiners one first fixes the
identifications of the endpoint Hilbert spaces. The product around the square is then central
when it acts trivially by conjugation on the full operator algebra. For irreducible observables
the center here is scalar, and the product of phases is therefore a number of unit modulus. A
phase assigned to each endpoint cancels around a closed square. For \eqref{eq:weylphase} that
number is exactly $e^{-ia_0 b_0/\hbar}$.

The square also indicates which type of observation could detect the discrepancy in the model.
Prepare a two-dimensional control in $(|0\rangle+|1\rangle)/\sqrt2$ and a normalized residual
state $f$. On branch zero act on $f$ by the identity, and on branch one act by $UVU^*V^*$ with
the same specified implementations. The final control has components $1$ and $e^{-ia_0
b_0/\hbar}$, up to their common normalization. Projecting it back on its initial state gives
probability
\[
 \frac{1+\cos(a_0 b_0/\hbar)}{2}.
\]
On the residual system alone, every expectation value is unchanged on both branches. This
calculation relies on a coherent implementation of the controlled operations; specifying just
the two automorphisms would not fix such an experiment. In a gravitational application one
would have to identify a preparation and a physically calibrated relative phase. A projective
phase of an isolated state has no separately measured value.

The step area $a_0 b_0$ can be arbitrarily small compared with $\hbar$ while leaving a
mathematically nontrivial phase. Repeating the loop $m$ times changes the phase to $e^{-im a_0
b_0/\hbar}$. The accumulation law follows from taking a product of the same central element. It
is a quantitative example of the difference between one local comparison of generators and a
comparison of transition amplitudes under repeated composition. It also explains why a
numerical verification of a single commuting square of algebra maps could consistently pass
while an amplitude calculation disagrees.

A reducible representation requires a more careful interpretation of the square. Two
implementers giving the same action on a represented algebra can differ by a unitary in its
commutant. For instance represent the same Weyl pair on $\mathcal H\otimes\CC^m$ by operators
$W(a,b)\otimes 1$. Every $1\otimes M$ commutes with them. An implementation of a repeated
operation can therefore act nontrivially on a multiplicity space even when it acts trivially on
every measured displacement. The scalar proof above used the full bounded algebra deliberately.
If a categorical description loses the multiplicity of a boundary preparation, identifying its
local observable algebra will not recover this part of the state comparison.

Chandrasekaran encounters an explicit endpoint phase when passing between reference actions in
his equation (4.3) and the discussion immediately following it \cite{nc}. There the change of
convention also changes the cap wavefunction, allowing the endpoint phases to cancel in an
amplitude. This mechanism calls for a rule on the states as well as on the algebra. A common
rephasing of endpoints and a nontrivial scalar left by a closed comparison have different
observational consequences, as the square calculation shows.

\section{Matching a positive boost charge}
\label{sec:boost}
There is a more direct contact with a gravitational calculation than an analogy between two
canonical pairs. In Chandrasekaran's Appendix D.2, equation (D.13) integrates the intermediate
rapidity of a shared normal frame. With his dimensionless area label $s=A/(4G_N\hbar)$, a
relative boost by $\Gamma$ acts on that sector with phase $e^{i\Gamma s/(2\pi)}$ \cite{nc}.
Matching two sectors with labels $s_L,s_R$ leaves a phase depending on the difference of their
rapidities. This formula tests three pieces of a prospective dictionary at once: its area
convention, its orientation convention and its normalization of the integration variable.

To isolate the Fourier calculation, take continuum labels and test functions with compact
support inside $(0,\infty)^2$. Suppose $\Gamma_L,\Gamma_R$ are fixed external readings. Give
the intermediate rapidity $\alpha$ a symmetric Gaussian window with scale $B>0$. The
corresponding kernel on the two label variables is
\begin{align}
 K_B(s_R,s_L)
 &=\int_{\RR}\frac{d\alpha}{4\pi^2}
       \exp\left(-\frac{\alpha^2}{2B^2}\right)
       \exp\left(\frac{i}{2\pi}
       [s_R(\Gamma_R-\alpha)+s_L(\alpha-\Gamma_L)]\right) \nonumber\\
 &=e^{i(s_R\Gamma_R-s_L\Gamma_L)/(2\pi)}
       \frac{B}{(2\pi)^{3/2}}
       \exp\left(-\frac{B^2(s_L-s_R)^2}{8\pi^2}\right).
 \label{eq:reg-match}
\end{align}
The second line is the elementary Fourier transform of a Gaussian. In the difference variable
the real factor is a normalized approximate identity of standard deviation $2\pi/B$. After
integration against the test function, its limit is evaluation on the diagonal. Compact support
away from zero permits extension by zero without any endpoint convention. The result is the
distribution
\[
 K_\infty(s_R,s_L)=
 e^{is_R(\Gamma_R-\Gamma_L)/(2\pi)}\,\delta(s_L-s_R).
\]
This reproduces the rapidity dependence and the normalization of the continuum Fourier identity
in (D.13). The Gaussian window here is a way of testing the distributional limit. It is
independent of a choice of Gaussian state for gravitational modes.

The result is sensitive to the normalization of the rapidity integral. If a construction
integrates a variable $\beta=c\alpha$ with $c>0$, it must transform both the exponent and the
measure. Keeping the physical phases fixed but integrating with $d\beta/(4\pi^2)$ in place of
$d\alpha/(4\pi^2)$ multiplies the entire kernel by $c$. The sharp matching locus looks
identical in both answers. Absolute amplitudes sewn with the two measures then disagree. Thus
the equality of constraints $s_L=s_R$ and equality of normalized sewing kernels are
independently testable statements.

\subsection{A comparison of spectral measures}
For a self-adjoint, nonnegative area operator $\widehat A$, a normalized state $\rho$ gives its
spectral probability measure $\mu_\rho$ on $[0,\infty)$. The boost probe has expectation
\begin{equation}
 \rho\!\left(e^{i\Gamma\widehat A/(8\pi G_N\hbar)}\right)
     =\int_{[0,\infty)} e^{i\Gamma A/(8\pi G_N\hbar)}\,d\mu_\rho(A).
 \label{eq:boost-marginal}
\end{equation}
Here the bounded functional calculus or a representation supporting it is understood. For a
vector state, $\mu_\rho(I)$ is the squared norm of the projection on the area interval $I$.
Agreement of \eqref{eq:boost-marginal} for every real $\Gamma$, with common constants and
rapidity convention, is equivalent to agreement of the probability measures. The forward
direction follows from uniqueness of the Fourier transform of a finite measure. To see why,
integrate the difference of their characteristic functions against the inverse Fourier
transform of a smooth compactly supported test function. Fourier inversion and absolute
integrability then give zero for the integral of every such test function against the measure
difference. These functions determine a finite measure.

For a sharp area eigenstate a single prescribed rapidity produces a periodic ambiguity. At
$\Gamma_0\ne0$, replacing $A$ by $A+2\pi (8\pi G_N\hbar)j/\Gamma_0$ for an integer $j$ leaves
that expectation unchanged whenever both labels are admissible. Two nonzero rapidities whose
ratio is irrational remove this exact ambiguity for a sharp label, since equality of both
phases would require the same nonzero area shift to be an integer period for both. For a
general distribution two probes still provide only two Fourier samples. Verifying them cannot
establish equality of all spectral probabilities. This distinction matters when an intended
comparison specifies a small set of preparation amplitudes instead of an entire state.

Even full equality of the area marginal does not settle the quantum state. To see this without
any domain subtleties, take $\mathcal H=L^2(\RR_+,dA)\otimes\CC^2$ with the area acting by
multiplication in the first factor. Fix a normalized packet $f(A)$ and take the two vectors
$f\otimes |0\rangle$ and $f\otimes |1\rangle$. Every boost and every bounded function of area
have identical expectation in the two states. The operator $1\otimes |0\rangle\langle0|$
distinguishes them. The second factor represents an undetected label in this calculation. The
gravitational proposal contains fluctuating radiative data within a fixed area in \S4.2, so a
comparison restricted to the boost subgroup would require a separate calculation on its
radiative observables.

\subsection{Continuous and atomic matching}
A Dirac delta on a continuum of labels and a projector on equal eigenvalues of a discrete
operator both express matching. They act on different test spaces. To compare them let $D$ be
the self-adjoint difference of the two area labels divided by $2\pi$ in a tensor product
representation. The functions of area on opposite factors strongly commute. A normalized finite
average of the corresponding one-parameter unitary group is
\[
 P_T=\frac1{2T}\int_{-T}^{T}e^{i\alpha D}\,d\alpha,\qquad
 P_T=\operatorname{sinc}(TD),\qquad
 \operatorname{sinc}(t)=
 \begin{cases}\sin(t)/t,&t\ne0,\\1,&t=0.\end{cases}
\]
For any vector its scalar spectral measure for $D$ is finite. The integrand multiplier has
absolute value at most one, equals one at zero, and tends to zero at every nonzero point.
Applying dominated convergence to the norm squared of the difference shows $P_T\to E_D(\{0\})$
strongly, where $E_D$ is the spectral resolution of $D$.

With pure point labels on each side this projection keeps vectors with equal labels along with
all their multiplicities. For discrete eigenvalues the list of allowed pairs and their
degeneracies can be computed before normalizing the physical trace. With an absolutely
continuous joint label measure on the plane, the diagonal has measure zero, hence
$E_D(\{0\})=0$. The normalized average then converges strongly to zero in the kinematical
Hilbert space. The unnormalized distributional pairing in \eqref{eq:reg-match} can nevertheless
be nonzero on smooth test functions. For example a smooth nonnegative bump positive at a
diagonal point has a positive diagonal integral when $\Gamma_L=\Gamma_R$.

In the latter case the diagonal data need a physical scalar product obtained after reduction,
just as $\eta$ did for a noncompact frame position. Substituting a symbol $\delta(0)$ for the
diagonal measure in a countable basis would conflate the two constructions. A proposed
identification of boundary Hilbert spaces must specify whether it preserves a spectral
subspace, a direct integral or the physical space completed after distributional matching. The
Fourier example decides which question to ask once the spectral measures are given. It does not
select a spectrum for the gravitational theory.

\subsection{The range of the area label}
The global range of a charge restricts the operations that can act on it. Suppose $\widehat A$
is a self-adjoint nonnegative operator on a nonzero Hilbert space. If for every $t\in\RR$ a
unitary $V(t)$ obeyed
\[
 V(t)\widehat A V(t)^*=\widehat A+t1
\]
as self-adjoint operators with their domains, then invariance of the spectrum under unitary
conjugation would give $\operatorname{spec}(\widehat A)=\operatorname{spec}(\widehat A)+t$ for
every $t$. The only nonempty subset of $\RR$ invariant under every translation is $\RR$ itself.
This contradicts nonnegativity. The assumption is a global unitary translation law. A
commutator identity on a chosen test domain alone does not imply this law, a distinction
studied in other semibounded canonical-pair problems by Galapon \cite{galapon}.

Accordingly the real Weyl pair of Section~\ref{sec:quantum} cannot be identified globally with
a positive area and a conjugate generating arbitrary additive shifts of that area. A local
canonical coordinate can still describe fluctuations in a regular chart, and the boost itself
acts by a phase at fixed area. These operations ask for substantially less than a global
two-sided shift of the area spectrum.

A simple alternative illustrates the importance of the range. On $L^2((0,\infty),dA/A)$ let
$\widehat A$ multiply by $A$ and define $(V(t)f)(A)=f(e^tA)$. The invariant measure $dA/A$
makes $V(t)$ unitary. On the self-adjoint multiplication domain it gives $V(t)\widehat A
V(t)^*=e^t\widehat A$. Hence conjugating a boost by a dilation rescales its rapidity:
\[
 V(t)\exp\bigl(i\Gamma\widehat A/(8\pi G_N\hbar)\bigr)V(t)^*
    =\exp\bigl(i e^t\Gamma\widehat A/(8\pi G_N\hbar)\bigr).
\]
Here zero belongs to the spectrum as a limit point and is preserved by dilation. In logarithmic
coordinates $x=\log(A/A_*)$ for a chosen positive reference area $A_*$, $V(t)$ acts by an
ordinary translation of $x$. This model gives a consistent pair of operations preserving
positivity. It is a comparison example; the gravitational admissibility of a dilation would
have to follow from the null-boundary action and its state space.

There is also a finite-dimensional obstruction to importing the polynomial Weyl algebra
unchanged into a projected size sector. For finite square matrices $X,Y$, cyclicity gives
$\operatorname{Tr}([X,Y])=0$. If $[X,Y]=i\hbar 1$ on a nonzero finite-dimensional space, the
right side has trace $i\hbar$ times its dimension, a contradiction. Projecting a canonical pair
defined on a larger space need not preserve its commutator. With a projection $P$ and bounded
operators $X,Y$ one has exactly
\[
 [PXP,PYP]=P[X,Y]P-PX(1-P)YP+PY(1-P)XP .
\]
Even for bounded inputs the two cross terms can carry the difference. For unbounded canonical
generators further common-domain conditions are needed before writing the same expression.
Chandrasekaran's program distinguishes a representation at the level of boundary observables
from bounds under a restriction on the size of a diamond. A purported comparison has to
identify the restriction and recompute what survives it. Declaring an algebraic Weyl relation
exact on a finite block would fail the trace test.

\subsection{Local smearing and the total area}
For a smooth regular cut, the integrated generator in Chandrasekaran's equation (3.43) sums a
charge aspect. Its value records the area of the cut but loses its angular distribution. In
\S4.2 the incoming data at fixed total area can fluctuate in the distribution of area over the
corner. Thus comparing a uniform boost on the whole cut and comparing separately labelled
probes on corner patches are physically different requests \cite{nc}.

A two-cell measurement makes the difference explicit with no assumption of full field
quantization. Suppose two labelled commuting positive observables $A_1,A_2$ take values $l$ and
$h$, with $0<l<h$. On states diagonal in the joint spectral basis the data are ordinary
probability distributions on four pairs. Consider preparations
\[
 \nu_+=\tfrac12\delta_{(l,l)}+\tfrac12\delta_{(h,h)},\qquad
 \nu_-=\tfrac12\delta_{(l,h)}+\tfrac12\delta_{(h,l)} .
\]
Each individual cell has the same probability for both outcomes in both preparations. All
single-cell boost characteristic functions consequently agree at all rapidities. Yet the total
area is either $2l$ or $2h$ in the first preparation and is always $l+h$ in the second. If $t$
is the calibrated coefficient of $A_1+A_2$ in the phase, the two expectations are
\[
 e^{it(l+h)}\cos\bigl(t(h-l)\bigr)
 \quad\hbox{and}\quad e^{it(l+h)}.
\]
Taking $t=\pi/(2(h-l))$ gives zero for the first expectation and unit modulus for the second.
The discrepancy is a cross correlation, and does not come from either one-cell spectral
measure.

There is also a converse failure. Sharp preparations at $(l,h)$ and $(h,l)$ give identical
total area and identical expectation for every uniform boost. An operation sensitive to the
first labelled cell resolves them. Equality of total boost data would be sufficient in a model
in which only the total area is an observable. Once a proposed class of decorated embeddings
includes a local patch smearing, the comparison needs additional state data.

In the two-cell algebra all joint information is accessible from
\[
 \chi(t_1,t_2)=
 \sum_{x,y\in\{l,h\}}\nu(x,y)\,e^{i(t_1x+t_2y)} .
\]
For instance choose $t_*=\pi/(h-l)$. The function $e^{it_*(x-l)}$ is $+1$ at $l$ and $-1$ at
$h$. The three expectations of the sign on the first cell, the sign on the second cell, and
their product determine the four probabilities along with normalization. This gives a finite
set of tests on the joint state in this particular spectrum. A physical corner would need
independently specified smearings, operators and a family of preparations before the same
reasoning could be applied.

A change of angular coordinates on a cut provides a kinematical consistency equation for such a
choice. Let a classical patch measurement on a surface $S$ be $A_S[f]=\int_S f\,d\mu_S$, with
$f$ a fixed scalar test function and $d\mu_S$ the area density. For an orientation-preserving
identification $\xi:S\to S'$ suppose $\xi^*d\mu_{S'}=d\mu_S$. Then
\[
 A_{S'}[f\circ\xi^{-1}]=A_S[f] .
\]
This follows from change of variables and is an equation on labelled measurements. A mere
assertion that two totals agree is its special case $f=1$. In a larger diamond the old corner
can cease to be part of the boundary, and a map $\xi$ of the required kind need not exist.
Extending the measurement into the larger gravitational algebra would then require an explicit
dressing or an equivalent protocol for recovering the patch. Assuming a map of angular
coordinates fixes that problem only for cuts already identified.

The finite example also clarifies a requirement on any entropy comparison. Let $M_d(\CC)$ act
on $\CC^d\otimes\CC^m$ as $a\mapsto a\otimes 1_m$. For any density matrices $\rho,\sigma$ on
the two factors, the state $\rho\otimes\sigma$ has expectation $\operatorname{Tr}(\rho a)$ on
every observable in the image. If $\sigma$ is the maximally mixed state, the eigenvalues of
$\rho\otimes\sigma$ are the eigenvalues of $\rho$ divided by $m$, each repeated $m$ times. Its
von Neumann entropy therefore exceeds that of $\rho$ by $\log m$. The representation has
changed while all expectations on the smaller algebra agree. A full entropy claim must
consequently identify which degrees of freedom and which trace enter the state count. In
particular equality of area distributions under all probes on a single cut is insufficient to
identify a gravitational entropy across different choices of observable algebra.

\section{Tests on null data}
The transverse geometry of a null sheet differs from the model of translations on $T^*V$.
Chandrasekaran's equations (3.39) and (3.42) exhibit shear and corner terms in the symplectic
potential and its exterior derivative. The area-dependent boost charge couples to variations
of the conformal metric and to changes of the null normals. In particular it would be
unjustified to infer a direct tensor factor for the corner charge from the independent pairs
in the present example. His Appendix D.1 also distinguishes first-class and second-class
constraints in its reduced measure. Equating every constraint with the generator of a group
action would lose this distinction.

For comparison, on one canonical plane the constraints $s=0$ and $n=0$ have $\{s,n\}=1$. Their
zero set is a point, and the matrix of mutual brackets is invertible. Reducing this
second-class pair by the Dirac bracket removes the plane \cite{dirac}. Tacking on a spectator pair $(r,k)$
leaves its bracket unchanged because the spectator functions commute with both constraints.
Treating $s,n$ as commuting generators of gauge translations on the common zero locus would
instead require a translation tangent to that locus. Neither translation is tangent there.
This finite example shows how the rank of the constraint bracket, a separate piece of input,
changes the meaning of a quotient.

A physical calculation can begin on a regular family of decorated diamonds. Choose a cut
without a generator crossing before the specified corner, choose two frame patches and specify
the observable carried by a common smearing at their interface. Fix which perturbations obey
the null constraints and boundary conditions. For these perturbations the kernel of the
pulled-back symplectic form determines the available degeneracies. A transformation with a
nonzero Hamiltonian surface charge is separately recorded whenever the boundary condition
permits it. At this step the algebraic comparison needs an actual map on the reduced variables
and on changes of frame.

At the quantum step one can ask whether a proposed transport preserves the chosen charge
sector, changes to another specified sector, or has an anomaly on a common domain.
Equations~\eqref{eq:ideal} and \eqref{eq:norm} are one formulation for normalizer
presentations. If one quantizes the reduced phase space directly, the analogous question is
whether the map respects its reduced brackets and quantum relations and can be implemented in
the chosen state class. A proof for boost observables alone would establish a comparison on
that subalgebra. The radiative fields and the couplings of the corner would remain to be
addressed.

A separate question concerns the source of the categorical observable. In the cited finite
reconstruction, relational weights give a metric on a graph \cite{rep}. To turn a collection
of probes into a copy of the corner algebra requires an independent prescription for
extracting the measured coordinates, relations, involution and gauge identifications. Reading
them from the null-cone algebra and declaring those objects to be the probes would make the
comparison tautological. An informative first observable would have a stated measurement on
each side and a test of equality after transporting it through one nonidentity arrow.

\section{Conclusion}
A common frame constraint leaves a relative reading and its conjugate charge measurable. This
conclusion holds in the degree-zero Poisson algebra and in the quantized normalizer quotient
of the translation model. The dependent tests used to describe the constraint persist in
negative Koszul degree. On a network, a quantum character inconsistent around a cycle kills
the full quotient even when it passes every test on a tree. These calculations put a concrete
condition on any proposed gravitational realization. It must carry the physical generators and
their quantum constraint levels consistently across changes of boundary frame.

The rank-changing family identifies an operator that disappears if one reduces before
specializing. At fixed rank the state comparison still detects inequivalent preparations
through their displacement widths, and two implementers can leave a phase around a square
that every conjugation on observables misses. At a null cut, one physical test is available
already. Two independent descriptions can compare the spectral distribution of the
area-generated boost, the normalized rapidity pairing, and the joint measurements carried by
identified corner patches. Agreement there would give a correspondence for specified
boundary data while leaving the radiative algebra to be determined.

\begin{thebibliography}{9}
\small
\bibitem{nc} V. Chandrasekaran, \emph{An algebraic quantization of gravity from the null cone} (2026). \href{https://arxiv.org/abs/2611.34924v1}{arXiv:2611.34924v1}.
\bibitem{df} W. Donnelly and L. Freidel, \emph{Local subsystems in gauge theory and gravity}, JHEP 09 (2016) 102. \href{https://arxiv.org/abs/1601.04744}{arXiv:1601.04744}.
\bibitem{long} M. Long, \emph{Quantum Gravity and Emergent Spacetime from Perspectival Structure}, YonedaAI Research Collective, 24 February 2026. \href{https://yonedaai.com/papers/pdf/quantum-gravity-emergent-spacetime.pdf}{Manuscript}.
\bibitem{rep} M. Long, \emph{A Representation-Theoretic Foundation for Quantum Gravity}, YonedaAI Research Collective, 7 July 2026. \href{https://yonedaai.com/papers/pdf/library/condensed-representation/part6-condensed-quantum-gravity.pdf}{Manuscript}.
\bibitem{derived} D. Arinkin, A. C\u{a}ld\u{a}raru and M. Hablicsek, \emph{Formality of derived intersections and the orbifold HKR isomorphism} (2014). \href{https://arxiv.org/abs/1412.5233}{arXiv:1412.5233}.
\bibitem{brst} M. Bordemann, H.-C. Herbig and S. Waldmann, \emph{BRST cohomology and phase space reduction in deformation quantization}, Commun. Math. Phys. 210 (2000), 107 to 144. \href{https://arxiv.org/abs/math/9901015v2}{arXiv:math/9901015v2}.
\bibitem{cahill} K. E. Cahill and R. J. Glauber, \emph{Ordered Expansions in Boson Amplitude Operators}, Phys. Rev. 177 (1969), 1857 to 1881. \href{https://doi.org/10.1103/PhysRev.177.1857}{doi:10.1103/PhysRev.177.1857}.
\bibitem{galapon} E. A. Galapon, \emph{Pauli's Theorem and Quantum Canonical Pairs: The Consistency of a Bounded, Self-Adjoint Time Operator Canonically Conjugate to a Hamiltonian with Non-empty Point Spectrum}, \href{https://arxiv.org/abs/quant-ph/9908033}{arXiv:quant-ph/9908033}.
\bibitem{dirac} P. A. M. Dirac, \emph{Generalized Hamiltonian Dynamics}, Canadian J. Math. 2 (1950), 129 to 148. \href{https://doi.org/10.4153/CJM-1950-012-1}{doi:10.4153/CJM-1950-012-1}.
\end{thebibliography}
\end{document}
